\documentclass[../../main.tex]{subfiles}
\begin{document}
\section{Stokes 定理}\label{section:ch16-stokess-theorem}

我们现在来到本章的中心结果。回忆一下，一个带边光滑流形 $M$ 的边界 $\partial M$ 自身是一个光滑流形，并带有一个诱导定向。

\begin{theorem}[Stokes 定理]\label{theorem:ch16-stokess-theorem}
设 $M$ 是一个定向的 $n$ 维光滑带边流形，$\omega$ 是 $M$ 上一个紧支集的光滑 $(n-1)$-形式。则
\begin{align}
\int_M\mathrm{d}\omega=\int_{\partial M}\omega.
\end{align}
\end{theorem}

\begin{figure}[ht]
\centering
\includegraphics[width=0.6\textwidth]{fig-16-9.png}
\caption{Stokes 定理}\label{figure:ch16-stokes}
\end{figure}

\begin{proof}这只需在定向坐标卡中证明，并用单位分解把结果拼接起来。在一个内部坐标卡中，$\omega$ 支集于 $\mathbb{R}^n$ 中一个紧立方体内部，$\partial M$ 上的积分为零，而 $\int_M\mathrm{d}\omega$ 是一个形如 $\int_{\mathbb{R}^n}\partial F/\partial x^i\,\mathrm{d}x^1\cdots\mathrm{d}x^n$ 的项之和；对每个 $i$ 作一次单变量积分即得零。在一个边界坐标卡中，像上面一样积分，除了一个端点落在 $x^n=0$ 上以外，其余各项都为零；这恰好给出 $\int_{\partial M}\omega$。
\end{proof}

\begin{remark}
Stokes 定理把几个经典的积分定理都作为特例囊括进来。当 $n=1$ 时，它就是微积分基本定理；当 $M\subseteq\mathbb{R}^2$ 是一个平面区域时，它就是 Green 公式；当 $M\subseteq\mathbb{R}^3$ 是一个区域时，它就是散度定理；而当 $M\subseteq\mathbb{R}^3$ 是一个曲面时，它就是经典的曲面 Stokes 公式。
\end{remark}

\begin{example}\label{example:ch16-circle-in-plane}
设 $M$ 是 $\mathbb{R}^2$ 中的单位圆盘 $\overline{\mathbb{B}}^2$，$\omega=(-y\,\mathrm{d}x+x\,\mathrm{d}y)/(x^2+y^2)$ 在除去原点外有定义。Stokes 定理不能直接用于整个圆盘，因为 $\omega$ 在原点附近不光滑。但在单位圆周上直接计算线积分给出 $2\pi$。
\end{example}

下面几个推论是 Stokes 定理的直接推论。

\begin{corollary}\label{corollary:ch16-closed-form-cycle}
设 $M$ 是一个定向光滑流形（不带边），$\omega$ 是 $M$ 上一个闭的光滑 $(n-1)$-形式。如果 $S\subseteq M$ 是一个紧的、定向的、光滑的 $(n-1)$ 维嵌入子流形（不带边），则 $\int_S\omega$ 只依赖于 $S$ 的同调类。特别地，若 $S$ 是某个紧带边子流形的边界，则 $\int_S\omega=0$。
\end{corollary}

\begin{corollary}\label{corollary:ch16-exact-form-boundary}
如果 $\omega=\mathrm{d}\eta$ 是 $M$ 上一个恰当形式，而 $S\subseteq M$ 是一个紧定向带边子流形，则
\begin{align*}
\int_{\partial S}\omega=0.
\end{align*}
\end{corollary}

\begin{corollary}\label{corollary:ch16-stokes-zero}
设 $M$ 是一个紧定向光滑流形（不带边）。如果 $\omega$ 是 $M$ 上一个恰当 $(n-1)$-形式，则
\begin{align*}
\int_M\omega=0.
\end{align*}
\end{corollary}

\begin{example}\label{example:ch16-circle-around-origin}
设 $\omega=(-y\,\mathrm{d}x+x\,\mathrm{d}y)/(x^2+y^2)$ 是 $\mathbb{R}^2\setminus\{0\}$ 上的光滑 $1$-形式，它是闭的但不是恰当的。绕原点一次的单位圆周 $S^1$ 满足 $\int_{S^1}\omega=2\pi$，而绕原点 $k$ 次的圆周给出 $2\pi k$。这表明闭形式的积分可以是任意整数倍的 $2\pi$。
\end{example}

\begin{theorem}[Green 公式]\label{theorem:ch16-greens-theorem}
设 $D\subseteq\mathbb{R}^2$ 是一个具有分段光滑边界的有界闭区域，带有标准的平面定向，$\partial D$ 带有诱导定向。如果 $P,Q$ 是 $D$ 的邻域上的光滑实值函数，则
\begin{align}
\int_D\left(\frac{\partial Q}{\partial x}-\frac{\partial P}{\partial y}\right)\mathrm{d}x\,\mathrm{d}y=\int_{\partial D}P\,\mathrm{d}x+Q\,\mathrm{d}y.
\end{align}
\end{theorem}

\begin{proof}令 $\omega=P\,\mathrm{d}x+Q\,\mathrm{d}y$，则 $\mathrm{d}\omega=(\partial Q/\partial x-\partial P/\partial y)\,\mathrm{d}x\wedge\mathrm{d}y$。这正是 $n=2$ 时的 Stokes 定理。
\end{proof}

\begin{example}\label{example:ch16-flux-through-circle}
设 $D$ 是 $\mathbb{R}^2$ 中以原点为心、半径为 $R$ 的闭圆盘。取 $P=-y$，$Q=x$，Green 公式给出
\begin{align*}
\int_{\partial D}x\,\mathrm{d}y-y\,\mathrm{d}x=\int_D 2\,\mathrm{d}x\,\mathrm{d}y=2\pi R^2,
\end{align*}
这正是圆盘面积的两倍，与圆周上的直接参数化计算一致。
\end{example}
\end{document}
